\chapter{Implementation and Evaluation}
\label{chap:impl}

The preceding chapter defined the GEO pipeline at the architectural level.
This chapter documents how that architecture was translated into a working system:
the technology choices made under real resource constraints, the non-trivial
engineering decisions that the design left open, the concrete failures encountered
during development and how they were resolved, and the experimental results
obtained on the Tunisian restaurant domain.

% ═══════════════════════════════════════════════════════════════════════════════
\section{Development Environment and Technology Stack}
\label{sec:devenv}
% ═══════════════════════════════════════════════════════════════════════════════

The system was implemented in \textbf{Python 3.11} on a standard Windows 10 machine,
without access to dedicated GPU hardware. This constraint is not incidental: it reflects
the reality of academic research conducted outside industrial infrastructure, and it
directly shaped every technology choice made throughout the project.

Three design imperatives governed the selection of components.
The first was \textit{cost accessibility}: all LLM inference had to operate within
free or low-cost API tiers, since sustained experimentation across hundreds of queries
would otherwise be financially prohibitive. This requirement alone motivated the design
of the dynamic \texttt{ModelRegistry}, which autonomously rotates across providers and
API keys to maximise available quota rather than relying on a single paid endpoint.
The second was \textit{architectural decoupling}: the interface between agents and
external tools had to be defined through a stable protocol rather than direct function
calls, so that tool implementations could be replaced or extended without modifying
agent logic. This constraint led to the adoption of the Model Context Protocol (MCP)
as the tool communication standard.
The third was \textit{resilient persistence}: given that a complete pipeline run spans
multiple API sessions subject to daily quota exhaustion, all intermediate results had
to be checkpointed incrementally to disk, allowing execution to resume at any stage
without reprocessing already-completed work.

These three imperatives are not independent. The need for free-tier compatibility
directly drives the registry's complexity, and the need for resumable execution is
only necessary because free-tier quotas force multi-session runs in the first place.
Table~\ref{tab:stack} summarises the technologies selected under these constraints.

\begin{table}[H]
\centering
\caption{Technology Stack}
\label{tab:stack}
\begin{tabular}{|p{3.2cm}|p{3.0cm}|p{6.8cm}|}
\hline
\textbf{Component} & \textbf{Technology} & \textbf{Role and Rationale} \\
\hline
Language & Python 3.11 & Native async support, dominant LLM tooling ecosystem \\
\hline
Orchestration & Custom ReAct + LangGraph state & Adaptive execution, decouples reasoning from tool calls \\
\hline
LLM Provider 1 & Groq API & Sub-second inference via dedicated GroqChip hardware \\
\hline
LLM Provider 2 & OpenRouter & Unified access to 1,610+ models, primary fallback pool \\
\hline
LLM Provider 3 & Mistral API (paid) & Reliable JSON output, independent quota pool for extraction \\
\hline
Model Management & Custom \texttt{ModelRegistry} & Per-model quota tracking and autonomous selection \\
\hline
Tool Protocol & FastMCP (stdio) & MCP-compliant servers, decoupled from agent code \\
\hline
Web Search & DuckDuckGo, SerpAPI & Complementary sources for entity web presence \\
\hline
Data Storage & pandas (CSV) & Incremental checkpointing, human-readable output \\
\hline
Fuzzy Matching & RapidFuzz & C-accelerated Levenshtein distance for entity deduplication \\
\hline
Encoding & UTF-8 BOM (\texttt{utf-8-sig}) & French/Arabic character support on Windows CP1252 \\
\hline
\end{tabular}
\end{table}

% ═══════════════════════════════════════════════════════════════════════════════
\section{System Implementation}
\label{sec:implementation}
% ═══════════════════════════════════════════════════════════════════════════════

The architecture defined in Chapter~\ref{chap:design} establishes the structural
blueprint of the pipeline but leaves a number of concrete implementation questions
open. This section documents the four decisions that most significantly shaped
the correctness, robustness, and validity of the final system.

% -----------------------------------------------------------------------
\subsection{Dynamic Model Management and the Slot Abstraction}
% -----------------------------------------------------------------------

A central requirement of the GEO analysis is that entity visibility be measured
\textit{comparably across capability tiers}: the same entity must be queried under
an 8B model, a 32B model, and a 70B model so that scale-dependent visibility can be
detected. Implementing this correctly requires a design decision that goes beyond
simple API calls  the \textit{identity} of the measurement instrument must remain
stable even when the underlying provider changes at runtime.

This is achieved through a \textbf{slot abstraction}: each model is represented by
a named slot (\texttt{llama-small}, \texttt{qwen-medium}, \texttt{llama-large})
rather than a specific model identifier. The slot name not the actual model ID is
what gets recorded in the output file. When Groq quota is exhausted, the registry
substitutes an OpenRouter equivalent within the same capability tier, but the slot
label remains unchanged. This ensures that per-tier GEO comparisons remain valid
across provider switches, which would otherwise introduce a confounding variable
between provider and model scale.

The \texttt{ModelRegistry} that implements this abstraction operates across three
providers and up to five API keys per provider, yielding a pool of 1,679 models at
initialisation. Model selection proceeds in four stages: (1) if the caller specifies
a preferred model and it is available, use it; (2) otherwise select the highest-tier
available model that fits the prompt's estimated token count; (3) if no model fits the
token budget but models are available, attempt with a warning; (4) if no model is
available at all, raise a \texttt{RuntimeError} with a breakdown of why each model is
blocked. Availability is determined by a per-model state machine that tracks three
independent conditions: TPD exhaustion (daily quota spent), TPM blocking (per-minute
limit active until a specific timestamp), and token capacity (learned from 413 errors).

A particularly important extension of this system is \textit{capacity learning}: when
a model returns an HTTP~413 error indicating that the request exceeds its token limit,
the registry records the reported capacity and uses it in subsequent calls to skip the
model for prompts that are known to exceed it. This eliminates repeated 413 failures
without requiring manual model configuration.

% -----------------------------------------------------------------------
\subsection{Two-Phase Entity Resolution}
% -----------------------------------------------------------------------

Raw LLM responses contain named entities in highly variable surface forms. The same
restaurant may appear as \textit{``Dar El Jeld''}, \textit{``dar el jeld''}, or
\textit{``Restaurant Dar El Jeld''} across different runs. Merging these variants into
a single canonical entity is not trivial: over-aggressive merging conflates distinct
establishments, while under-merging inflates the entity count and fragments mention
statistics.

The pipeline resolves this through a two-phase approach. \textbf{Phase B} applies
fuzzy string matching using \texttt{RapidFuzz}'s \texttt{token\_sort\_ratio}, which
handles word-order variance (\textit{``el ksar''} vs \textit{``ksar el''}) and is
C-accelerated for performance. Two thresholds govern the outcome:
a score above 88 triggers automatic merge; a score between 75 and 87 flags the pair
as ambiguous and defers it to Phase C.

\textbf{Phase C} submits ambiguous pairs to an LLM arbitrator with a conservative
default of DIFFERENT: two names are merged only if the evidence is high-confidence,
i.e., the difference is a pure spelling variant, article prepend, or transliteration.
Structurally distinct names even with high fuzzy similarity are kept separate.
The same LLM call also validates entity legitimacy: it checks whether the entity is
a real food-service establishment in Tunisia or North Africa, and rejects LLM
hallucinations such as Parisian restaurants appearing in responses about Tunisian
cuisine.

% -----------------------------------------------------------------------
\subsection{GEO Visibility Scoring}
% -----------------------------------------------------------------------

Once entities are cleaned and canonicalised, Agent~1 computes a set of
GEO-specific features for each entity. The core metric is the \textbf{mention rate}:
the fraction of all responses (across all prompts, slots, and runs) in which the
entity appears. However, mention rate alone is insufficient; an entity mentioned
once across 450 responses and one mentioned consistently across 10 responses may
have the same raw rate but very different GEO values.

The \textbf{stability score} addresses this by jointly penalising low mention
frequency and high variance in ranking position:

\begin{equation}
\text{stability} \;=\; \text{mention\_rate} \;\times\; \frac{1}{1 + \sigma_{\text{rank}}}
\label{eq:stability}
\end{equation}

Where $\sigma_{\text{rank}}$ is the standard deviation of ranking positions across
all mentions of the entity. An entity mentioned in every response at position 1
achieves a stability score of 1.0; one mentioned rarely and at inconsistent positions
approaches 0. Entities are classified as \textit{stable} (score $\geq 0.75$),
\textit{variable} ($0.40 \leq$ score $< 0.75$), or \textit{unstable} (score $< 0.40$).

Three additional features capture complementary dimensions of visibility:
\texttt{mention\_prominence} (fraction of mentions appearing in the top-3 ranking positions),
\texttt{cross\_model\_rate} (fraction of distinct model tiers that mention the entity),
and \texttt{co\_mention\_rate} (how frequently the entity co-appears with other top
entities in the same response). Together these features constitute the GEO profile
of each entity, enabling a multidimensional ranking that captures both frequency
and positional consistency.

% -----------------------------------------------------------------------
\subsection{Fault-Tolerant Execution and Rate Limit Management}
% -----------------------------------------------------------------------

Free-tier APIs impose two orthogonal quota types that require fundamentally opposite
recovery strategies. A \textit{Tokens-Per-Minute} (TPM) limit is transient: the
affected model recovers within seconds and should be retried after a short wait.
A \textit{Tokens-Per-Day} (TPD) limit is session-terminal: the model is unusable
until the UTC midnight quota reset, and retrying it wastes both time and the remaining
attempt budget allocated to the current entity. Misclassifying these two error types
has severe consequences in either direction: treating a TPD error as TPM causes the
pipeline to stall indefinitely; treating a TPM error as TPD permanently discards a
usable model from the session.

The \texttt{SupervisorAgent} handles this classification through priority-ordered
keyword matching that evaluates TPD signals \textit{before} TPM signals, since a
TPD error message may also contain rate-related vocabulary. A particularly subtle
case encountered during development is OpenRouter's \texttt{free-models-per-day}
error string, which contains the word ``per'' in a hyphenated compound that naive
substring matching fails to recognise. This exact string is included explicitly in
the TPD detection list to prevent misclassification.
Table~\ref{tab:supervisor} summarises the full classification logic.

\begin{table}[H]
\centering
\caption{Supervisor Error Classification Rules}
\label{tab:supervisor}
\begin{tabular}{|p{4.5cm}|p{3.0cm}|p{5.5cm}|}
\hline
\textbf{Error Signal} & \textbf{Action} & \textbf{Notes} \\
\hline
HTTP 429 + ``per day'' / ``free-models-per-day'' / HTTP 402 & \texttt{switch\_model} & Mark TPD-exhausted, persist timestamp to JSON \\
\hline
HTTP 429 + ``per minute'' / ``RPM'' & \texttt{wait\_retry} & Extract wait time from error message \\
\hline
HTTP 413 + token limit & \texttt{switch\_model} & Registry learns model capacity, skips for oversized prompts \\
\hline
HTTP 500 / 502 / 503 & \texttt{retry} & Transient server error, immediate retry \\
\hline
5+ consecutive RPM errors & \texttt{switch\_model} & Escape RPM loop, model effectively unusable \\
\hline
attempt $\geq 4$ & \texttt{skip} & Write error row, continue to next entity \\
\hline
\end{tabular}
\end{table}

Resumable execution is enforced at every pipeline stage through a \texttt{done\_keys}
set loaded at startup from the existing output file. For Agent~2 specifically, the
filter requires both \texttt{confidence > 0} and \texttt{error is null}: rows written
during failed attempts are excluded from the done set and retried on the next run.
This distinction is critical without it, quota exhaustion during Agent~2 execution
would permanently discard web data for every entity being processed at the moment of
failure.

% ═══════════════════════════════════════════════════════════════════════════════
\section{Engineering Challenges and Solutions}
\label{sec:challenges}
% ═══════════════════════════════════════════════════════════════════════════════

The implementation decisions described above were not arrived at deductively.
They emerged from a development process that surfaced a set of concrete failures,
each of which exposed a gap between an assumption encoded in the software and the
actual behaviour of real API providers and operating-system environments.
Table~\ref{tab:challenges} documents the principal failures, their root causes,
and the solutions applied.

\begin{table}[H]
\centering
\caption{Engineering Challenges and Implemented Solutions}
\label{tab:challenges}
\begin{tabular}{|p{4cm}|p{4.6cm}|p{4.6cm}|}
\hline
\textbf{Challenge} & \textbf{Root Cause} & \textbf{Solution} \\
\hline
Pipeline stalls on TPD exhaustion & No persistent per-model daily quota state & Registry persists exhaustion timestamps to JSON; survives restarts \\
\hline
Infinite RPM retry loop & Attempt counter decremented on wait-retry path, keeping it perpetually at 1 & Removed decrement; each wait consumes the attempt budget \\
\hline
\texttt{free-models-per-day} misclassified as TPM & Hyphenated OpenRouter string not matched by ``per day'' substring & Exact string added to TPD list; TPD checked before TPM \\
\hline
First CSV column name corrupted & Missing \texttt{utf-8-sig} codec on Windows CP1252 prepends BOM bytes to first column & Consistent \texttt{utf-8-sig} encoding on all read and write operations \\
\hline
Agent 0 overwrites French prompt set & LLM defaults to English on regeneration; no file protection & \texttt{agent0\_run()} never overwrites an existing prompt CSV \\
\hline
Error rows not retried on resume & \texttt{\_load\_done()} counted all rows including failed attempts & Filter: \texttt{confidence > 0} AND \texttt{error} is null \\
\hline
Agent 2 selecting non-tool-capable models & \texttt{registry.select()} returns any available fallback when preferred exhausted & Guard rejects returned model if not in \texttt{\_TOOL\_CAPABLE} whitelist \\
\hline
8B models reject Agent 2 system prompt & Agent 2 system prompt exceeds 6,000-token context of small models & Registry learns capacity from 413 errors; skips model for oversized prompts \\
\hline
\end{tabular}
\end{table}

Several of these failures share a common pattern: a silent mismatch between a
software assumption and the behaviour of an external system. The BOM bug assumed
that Python's default codec handles all Unicode correctly on Windows.
The TPD misclassification assumed that provider error strings follow a consistent
format. The empty-file bug assumed that file existence implies valid content.
In each case, the failure was invisible until triggered under live API conditions,
confirming that empirical testing under real quota and encoding constraints is
irreplaceable as a debugging method for multi-provider LLM pipelines.

% ═══════════════════════════════════════════════════════════════════════════════
\section{Experimental Setup}
\label{sec:setup}
% ═══════════════════════════════════════════════════════════════════════════════

\subsection{Domain and Motivation}

The experiment targets \textbf{Tunisian restaurants} food-service establishments
across all Tunisian cities and price segments. This domain was selected for three reasons.
First, it represents a local market with strongly heterogeneous international visibility:
a small number of establishments are documented internationally while the majority remain
unknown outside Tunisia, making the GEO visibility gap maximally observable.
Second, the French-Arabic bilingual character of Tunisian online content provides a
realistic test of cross-lingual entity normalisation.
Third, published GEO research has focused almost exclusively on large English-language
markets~\cite{aggarwal2023geo}; this study contributes the first analysis of a
North African, Francophone context.

\subsection{Prompt Set}

The prompt set consists of 30 French-language queries covering six intent categories
(Table~\ref{tab:prompts}), with five variants per intent.
French was chosen to simulate the dominant digital interaction register of Tunisian users.
Every prompt names a specific Tunisian city and uses superlative or ranking formulations
(\textit{``Quel est le meilleur...''}, \textit{``Cite-moi les 3 plus connus...''})
to ensure that responding LLMs produce named entity lists rather than generic descriptions.
The set was validated by Agent~0's self-reflection mechanism, which confirmed 100\%
brand-eliciting coverage and six balanced intent categories spanning informational,
locational, and contextual query types.

\begin{table}[H]
\centering
\caption{Prompt Set: Intent Categories and Examples}
\label{tab:prompts}
\begin{tabular}{|p{4cm}|p{4.7cm}|p{5.2cm}|}
\hline
\textbf{Intent} & \textbf{Description} & \textbf{Example Prompt} \\
\hline
\texttt{top\_recommendation} & Best restaurants by overall reputation & \textit{``Quel est le meilleur restaurant tunisien à Tunis ?''} \\
\hline
\texttt{cuisine\_authentique} & Traditional dish recommendations & \textit{``Quel restaurant est réputé pour son couscous à Tunis ?''} \\
\hline
\texttt{localisation} & Restaurants in a specific district & \textit{``Quels restaurants recommandes-tu dans le quartier La Marsa ?''} \\
\hline
\texttt{occasion\_speciale} & Context-specific venue recommendations & \textit{``Quel restaurant pour un dîner romantique à Tunis ?''} \\
\hline
\texttt{comparatif} & Explicit ranking and comparison & \textit{``Classe les 5 meilleurs restaurants tunisiens à Tunis.''} \\
\hline
\texttt{avis\_touristes} & Tourist-oriented queries & \textit{``Je suis touriste à Tunis, quel restaurant me recommandes-tu ?''} \\
\hline
\end{tabular}
\end{table}

\subsection{Query Configuration}

Each prompt is evaluated across three model slots
(\texttt{llama-small} 8B, \texttt{qwen-medium} 32B, \texttt{llama-large} 70B)
for $R = 5$ independent runs each, yielding a planned total of:

\[
N \;=\; P \times S \times R \;=\; 30 \times 3 \times 5 \;=\; 450 \;\text{LLM queries}
\]

The three-tier model design allows the study to quantify whether entity visibility
is consistent across model scale, or whether certain entities are known only to
larger models. The five-run repetition supports mention stability estimation:
an entity appearing in all five runs of a given (prompt, slot) pair exhibits high
stability, while one appearing in a single run signals low-reliability visibility.

% ═══════════════════════════════════════════════════════════════════════════════
\section{Results}
\label{sec:results}
% ═══════════════════════════════════════════════════════════════════════════════



\subsection{Agent 0 Prompt Set Validation}


\subsection{Agent 1 Response Collection and Entity Extraction}


\subsection{Agent 2 Web Presence Research}



\subsection{Agent 3  GEO Scoring and Visibility Ranking}



% ═══════════════════════════════════════════════════════════════════════════════
\section{Discussion}
\label{sec:discussion}
% ═══════════════════════════════════════════════════════════════════════════════

\subsection{Architecture Validation}


\subsection{Limitations}


% ═══════════════════════════════════════════════════════════════════════════════
\section{Conclusion}
% ═══════════════════════════════════════════════════════════════════════════════
